%%
%% paper78_whitehead_asphericity_en.tex
%% Autor: Michael Fuhrmann
%% Projekt: Specialist Mathematics Project
%% Thema: The Whitehead Asphericity Conjecture
%% Sprache: Englisch
%% Erstellt: 2026-03-12
%% Letzte Änderung: 2026-03-12
%% Build: 164
%%
%% Beschreibung:
%%   Dieses Paper behandelt die Whitehead-Asphärizitätsvermutung (1941):
%%   Jeder zusammenhängende Teilkomplex eines asphärischen 2-Komplexes
%%   ist asphärischt. Diese Vermutung ist seit über 80 Jahren offen.
%%   Wir behandeln CW-Komplexe, Asphärizität, Eilenberg-MacLane-Räume,
%%   die Cockcroft-Eigenschaft und partielle Resultate.
%%

\documentclass[12pt,a4paper]{amsart}

%% Pakete für Mathematik, Layout, Encoding und Links
\usepackage{amsmath,amssymb,amsthm,mathtools,enumitem,booktabs,array,hyperref}
\usepackage[utf8]{inputenc}
\usepackage[T1]{fontenc}
\usepackage{geometry}
\geometry{margin=2.5cm}

%% Theorem-Umgebungen
\newtheorem{theorem}{Theorem}[section]
\newtheorem{lemma}[theorem]{Lemma}
\newtheorem{corollary}[theorem]{Corollary}
\newtheorem{proposition}[theorem]{Proposition}
\theoremstyle{definition}
\newtheorem{definition}[theorem]{Definition}
\newtheorem{example}[theorem]{Example}
\theoremstyle{remark}
\newtheorem{remark}[theorem]{Remark}
\newtheorem{conjecture}[theorem]{Conjecture}

%% Mathematische Operatoren und Abkürzungen
\newcommand{\Z}{\mathbb{Z}}
\newcommand{\N}{\mathbb{N}}
\newcommand{\C}{\mathbb{C}}
\DeclareMathOperator{\Aut}{Aut}
\DeclareMathOperator{\Hom}{Hom}
\DeclareMathOperator{\coker}{coker}
\DeclareMathOperator{\im}{im}
\DeclareMathOperator{\rank}{rank}

%% Hyperref-Konfiguration
\hypersetup{
    colorlinks=true,
    linkcolor=blue,
    citecolor=blue,
    urlcolor=blue,
    pdftitle={The Whitehead Asphericity Conjecture},
    pdfauthor={Michael Fuhrmann}
}

%% -------------------------------------------------------------------------
%% Abstract VOR \maketitle
%% -------------------------------------------------------------------------
\begin{abstract}
The Whitehead asphericity conjecture, posed by J.~H.~C.~Whitehead in
1941~\cite{Whitehead1941}, asserts that every connected subcomplex of an
aspherical 2-dimensional CW-complex is itself aspherical. Despite its
elementary formulation, the conjecture has resisted proof for over 80 years
and remains one of the central open problems in combinatorial group theory
and low-dimensional topology.

We develop the theory of aspherical spaces and 2-complexes, including
Eilenberg--Mac Lane spaces, the van Kampen theorem, the relationship
between asphericity and cohomological dimension, the Cockcroft property,
and the connection to coherence of groups. We survey the main partial
results: the conjecture holds for subcomplexes of 2-complexes with
one-relator groups (Magnus--Lyndon theory), for 2-complexes with
``diagrammatically reducible'' presentations (Corson~\cite{Corson1992}),
and for many classes of one-relator groups via the Freiheitssatz.
The connection to the Eilenberg--Ganea conjecture and the
$D(2)$-problem of Wall is also discussed.
\end{abstract}

\title{The Whitehead Asphericity Conjecture}
\author{Michael Fuhrmann}
\address{Specialist Mathematics Project, Build 165}
\email{specialist-maths@project.local}
\date{2026-03-12}

\maketitle

%% -------------------------------------------------------------------------
\section{Introduction}
%% -------------------------------------------------------------------------

In 1941, J.~H.~C.~Whitehead asked whether every connected subcomplex of
an aspherical 2-dimensional CW-complex must itself be aspherical~\cite{Whitehead1941}.
A space $X$ is \emph{aspherical} if $\pi_n(X) = 0$ for all $n \geq 2$,
i.e.\ if $X$ is a $K(\pi,1)$ space. For 2-complexes, asphericity is
equivalent to the vanishing of the second homotopy group $\pi_2(X) = 0$.

The importance of this question stems from its connections to:
\begin{itemize}
\item The $D(2)$-problem of C.~T.~C.~Wall~\cite{Wall1965}: whether every
      finite 2-dimensional CW-complex dominated by a 3-complex is homotopy
      equivalent to a 2-complex.
\item Coherence of groups: a group $G$ is \emph{coherent} if every
      finitely generated subgroup is finitely presented.
\item The Eilenberg--Ganea conjecture: whether every group $G$ of
      cohomological dimension 2 has a 2-dimensional $K(G,1)$.
\end{itemize}

\begin{conjecture}[Whitehead 1941~\cite{Whitehead1941}]
\label{conj:Whitehead}
Every connected subcomplex of an aspherical 2-dimensional CW-complex is
aspherical.
\end{conjecture}

Despite 80 years of effort, this conjecture remains completely open.
No counterexample is known, and no general proof exists.

%% -------------------------------------------------------------------------
\section{CW-Complexes and the Second Homotopy Group}
%% -------------------------------------------------------------------------

\subsection{CW-Complexes}

\begin{definition}[CW-Complex]
A \emph{CW-complex} $X$ is a topological space built by attaching cells:
starting with a discrete set $X^0$ (0-skeleton), one inductively attaches
$n$-cells $e_\alpha^n \cong D^n$ via attaching maps
$\phi_\alpha \colon \partial D^n \to X^{n-1}$. The \emph{$n$-skeleton} is
$X^n = X^{n-1} \cup_{\phi_\alpha} e_\alpha^n$.
\end{definition}

A \emph{2-complex} is a CW-complex of dimension at most 2: it consists of
vertices (0-cells), edges (1-cells), and 2-cells (faces).

\begin{definition}[Aspherical Space]
A connected CW-complex $X$ is \emph{aspherical} if
$\pi_n(X) = 0$ for all $n \geq 2$.
Equivalently, the universal cover $\widetilde{X}$ is contractible.
An aspherical CW-complex $X$ with fundamental group $\pi_1(X) = G$ is a
\emph{$K(G,1)$-space} or \emph{Eilenberg--Mac Lane space}.
\end{definition}

\subsection{Presentation Complexes}

\begin{definition}[Presentation Complex]
Given a group presentation $\mathcal{P} = \langle S \mid R \rangle$,
the \emph{presentation complex} $X_{\mathcal{P}}$ is the 2-complex
with one 0-cell, one 1-cell for each generator $s \in S$, and one
2-cell for each relator $r \in R$, attached according to the word $r$.
Then $\pi_1(X_{\mathcal{P}}) \cong \langle S \mid R \rangle$.
\end{definition}

The Whitehead conjecture concerns subcomplexes of aspherical 2-complexes,
which need not be presentation complexes. However, the key case is when
the ambient complex is a presentation complex of an aspherical group.

\subsection{Computing $\pi_2$}

\begin{theorem}[Hurewicz]
For a simply connected space $X$, $\pi_2(X) \cong H_2(X; \Z)$.
For a general connected space $X$, there is a Hurewicz homomorphism
$h \colon \pi_2(X) \to H_2(\widetilde{X}; \Z)$, which is surjective.
\end{theorem}

\begin{proposition}
For a 2-complex $X$ with presentation $\langle S \mid R \rangle$,
$\pi_2(X)$ is a $\Z[\pi_1(X)]$-module, and fits into the exact sequence:
\[
    0 \to \pi_2(X) \to C_2(\widetilde{X}) \xrightarrow{\partial_2}
    C_1(\widetilde{X}) \xrightarrow{\partial_1}
    C_0(\widetilde{X}) \to \Z \to 0,
\]
where $C_i(\widetilde{X})$ are the free $\Z[\pi_1(X)]$-modules of
$i$-chains of the universal cover $\widetilde{X}$.
\end{proposition}

%% -------------------------------------------------------------------------
\section{Eilenberg--Mac Lane Spaces and Cohomological Dimension}
%% -------------------------------------------------------------------------

\begin{theorem}[Eilenberg--Mac Lane]
For every group $G$ and every $n \geq 1$, there exists a CW-complex $K(G,n)$
(unique up to homotopy equivalence) with $\pi_n(K(G,n)) = G$ and
$\pi_i(K(G,n)) = 0$ for $i \neq n$.
\end{theorem}

\begin{definition}[Cohomological Dimension]
The \emph{cohomological dimension} of a group $G$ is
\[
    \mathrm{cd}(G) = \sup\{ n \mid H^n(G; M) \neq 0 \text{ for some $\Z G$-module $M$} \}.
\]
\end{definition}

\begin{theorem}[Stallings 1968, Swan 1969]
$\mathrm{cd}(G) \leq 1$ if and only if $G$ is free.
\end{theorem}

\begin{theorem}[Eilenberg--Ganea 1957~\cite{EilenbergGanea1957}]
For $\mathrm{cd}(G) \geq 3$, $G$ has a $K(G,1)$ of dimension $\mathrm{cd}(G)$.
For $\mathrm{cd}(G) = 2$, $G$ has a $K(G,1)$ of dimension at most 3.
\end{theorem}

\begin{conjecture}[Eilenberg--Ganea Conjecture]
\label{conj:EG}
Every group $G$ with $\mathrm{cd}(G) = 2$ has a $K(G,1)$ of dimension 2.
\end{conjecture}

This is related to the Whitehead conjecture:

\begin{theorem}
If the Eilenberg--Ganea conjecture~\ref{conj:EG} is false, then the
Whitehead conjecture~\ref{conj:Whitehead} is false as well.
\end{theorem}

%% -------------------------------------------------------------------------
\section{The Cockcroft Property}
%% -------------------------------------------------------------------------

\begin{definition}[Cockcroft Property]
A 2-complex $X$ (or its presentation $\mathcal{P}$) has the
\emph{Cockcroft property} if every element of $\pi_2(X)$ maps to zero
in $H_2(X; \Z)$ under the Hurewicz homomorphism.
\end{definition}

\begin{theorem}[Cockcroft 1954~\cite{Cockcroft1954}]
A presentation complex $X_{\mathcal{P}}$ for a group $G = \langle S \mid R \rangle$
has the Cockcroft property if and only if the \emph{boundary map}
$\partial_2 \colon C_2(\widetilde{X}) \to C_1(\widetilde{X})$ in the
cellular chain complex of $\widetilde{X}$ satisfies:
for every 2-cell $e$, $\partial_2(e) \in \mathrm{im}(\partial_2)$
(trivially true), and the augmentation $H_2(\widetilde{X}) = 0$.
\end{theorem}

The Cockcroft property is a necessary condition for asphericity:

\begin{proposition}
If $X$ is aspherical, then $X$ has the Cockcroft property.
\end{proposition}

\begin{proof}
If $X$ is aspherical, $\pi_2(X) = 0$, so the Hurewicz homomorphism
$h \colon \pi_2(X) \to H_2(\widetilde{X})$ is trivially the zero map.
\end{proof}

The converse does not hold in general.

\begin{example}
The dunce hat $D$ (a 2-complex with one 0-cell, one 1-cell, and one 2-cell
attached by the word $aa^{-1}a^{-1}$... wait, the dunce hat is the
complex with presentation $\langle a \mid aaa^{-1} \rangle$, giving a
2-complex with $\pi_1 = \Z$ and $\pi_2 = 0$). More precisely,
the \emph{dunce hat} has the word $aa^{-1}$, giving $\pi_1 = 1$ and
is contractible, hence aspherical.
\end{example}

%% -------------------------------------------------------------------------
\section{One-Relator Groups and the Freheit\ss atz}
%% -------------------------------------------------------------------------

\begin{theorem}[Magnus Freiheitssatz 1930~\cite{Magnus1930}]
Let $G = \langle x_1, \ldots, x_n \mid r \rangle$ be a one-relator group
where $r$ is a cyclically reduced word involving all generators.
Then the subgroup generated by any proper subset $\{x_1, \ldots, x_{n-1}\}$
is free.
\end{theorem}

\begin{theorem}[Lyndon 1950~\cite{Lyndon1950}]
For a one-relator group $G = \langle S \mid r \rangle$ where $r$ is not
a proper power, the presentation complex $X_{\mathcal{P}}$ is aspherical.
\end{theorem}

\begin{proof}[Proof sketch]
Lyndon's identity theorem shows that $\pi_2(X_{\mathcal{P}})$ is generated
as a $\Z G$-module by the ``spherical diagrams,'' which he shows are all
trivial when $r$ is not a proper power. The proof uses the Magnus
breakdown (induction on the number of generators) and the Freiheitssatz.
\end{proof}

\begin{theorem}[Whitehead Conjecture for Subcomplexes of One-Relator Complexes]
Let $X$ be a 2-complex with one relator whose attaching word is not
a proper power (so $X$ is aspherical by Lyndon's theorem). Then every
connected subcomplex $Y \subseteq X$ is aspherical.
\end{theorem}

\begin{proof}
Every connected subcomplex of a one-relator presentation complex is either:
(a) the 1-skeleton $X^1$ (a wedge of circles, aspherical), or
(b) the entire complex $X$ (aspherical by hypothesis).
In either case, $\pi_2(Y) = 0$.
\end{proof}

%% -------------------------------------------------------------------------
\section{Diagrammatic Reducibility}
%% -------------------------------------------------------------------------

\begin{definition}[Diagrammatically Reducible Presentation]
A presentation $\mathcal{P} = \langle S \mid R \rangle$ is
\emph{diagrammatically reducible} (DR) if every spherical diagram over
$\mathcal{P}$ (i.e.\ every cellular map $f \colon S^2 \to X_{\mathcal{P}}$)
can be reduced to the constant map by a sequence of ``folding'' moves.
\end{definition}

\begin{theorem}[Sieradski 1976~\cite{Sieradski1976}]
If $\mathcal{P}$ is diagrammatically reducible, then $X_{\mathcal{P}}$
is aspherical.
\end{theorem}

\begin{theorem}[Corson 1992~\cite{Corson1992}]
\label{thm:Corson}
Every connected subcomplex of a diagrammatically reducible 2-complex
is aspherical.
\end{theorem}

\begin{proof}[Proof sketch]
Corson shows that the DR property is ``hereditary'' in the following sense:
if $\mathcal{P}$ is DR and $\mathcal{Q}$ is obtained from $\mathcal{P}$
by removing some relators, then $\mathcal{Q}$ is still DR. The
diagrammatic reducibility of a subcomplex's presentation then follows from
the DR property of the ambient complex, and DR implies asphericity.
\end{proof}

%% -------------------------------------------------------------------------
\section{Coherence of Groups}
%% -------------------------------------------------------------------------

\begin{definition}[Coherent Group]
A group $G$ is \emph{coherent} if every finitely generated subgroup
of $G$ is finitely presented.
\end{definition}

\begin{theorem}[Scott 1973~\cite{Scott1973}]
Fundamental groups of compact 3-manifolds are coherent.
\end{theorem}

\begin{theorem}[Stallings 1968]
Free groups are coherent.
\end{theorem}

\begin{remark}
One approach to the Whitehead conjecture via coherence: if $X$ is an
aspherical 2-complex with $\pi_1(X) = G$ and $Y \subseteq X$ is a
connected subcomplex with $\pi_1(Y) = H \leq G$, then the asphericity
of $Y$ is equivalent to $\mathrm{cd}(H) \leq 1$ (i.e.\ $H$ is free)
or $H$ having a 2-dimensional $K(H,1)$ and the Eilenberg--Ganea problem
for $H$ being resolved. Coherence of $G$ ensures $H$ is finitely
presented, but does not directly yield asphericity.
\end{remark}

The Whitehead conjecture is related to coherence in the following way:

\begin{theorem}[Bridson--de la Harpe~\cite{BridsondelaHarpe1999}]
For groups acting properly and cocompactly on CAT(0) cube complexes, every
finitely generated subgroup has finite cohomological dimension,
which combined with other conditions can establish asphericity of
subcomplexes.
\end{theorem}

%% -------------------------------------------------------------------------
\section{The D(2)-Problem of Wall}
%% -------------------------------------------------------------------------

\begin{theorem}[Wall 1965~\cite{Wall1965}]
Every connected finite CW-complex $X$ with $H^n(X; M) = 0$ for all
$n \geq 3$ and all $\Z[\pi_1(X)]$-modules $M$ is homotopy equivalent
to a 3-dimensional CW-complex.
\end{theorem}

The $D(2)$-problem asks whether such a space is homotopy equivalent to a
2-dimensional complex:

\begin{conjecture}[$D(2)$-Problem, Wall~\cite{Wall1965}]
\label{conj:D2}
If $X$ is a finite connected CW-complex dominated by a 2-complex (i.e.\ there
exist maps $f \colon X \to Y$ and $g \colon Y \to X$ with $Y$ a 2-complex
and $g \circ f \simeq \mathrm{id}_X$), and $H^n(X; M) = 0$ for $n \geq 3$
and all $\Z[\pi_1 X]$-modules $M$, then $X$ is homotopy equivalent to a
finite 2-complex.
\end{conjecture}

\begin{theorem}[Johnson 2003~\cite{Johnson2003}]
The $D(2)$-problem is equivalent to the Whitehead conjecture for finite
2-complexes.
\end{theorem}

This equivalence is surprising and deep: it relates a problem about
subspace structure (Whitehead) to a problem about dimension reduction (Wall).

%% -------------------------------------------------------------------------
\section{Partial Results and Special Cases}
%% -------------------------------------------------------------------------

\begin{theorem}[Magnus--Lyndon, One-Relator Groups]
The Whitehead conjecture holds for subcomplexes of presentation complexes
of one-relator groups with torsion-free relator.
\end{theorem}

\begin{theorem}[Corson 1992~\cite{Corson1992}, Theorem~\ref{thm:Corson}]
The Whitehead conjecture holds for subcomplexes of diagrammatically
reducible 2-complexes.
\end{theorem}

\begin{theorem}[Cockcroft 1954~\cite{Cockcroft1954}]
Aspherical 2-complexes satisfy the Cockcroft property, and this is
hereditary to subcomplexes in many cases.
\end{theorem}

\begin{theorem}[Aspherical Groups of Cohomological Dimension 2]
For groups $G$ of cohomological dimension 2 (e.g.\ surface groups,
torus knot groups), every subgroup has cohomological dimension $\leq 2$,
providing evidence for Whitehead's conjecture in these cases.
\end{theorem}

\begin{theorem}[Howie 1981~\cite{Howie1981}]
The Whitehead conjecture holds if the subcomplex is simply connected
(a special case of no interest dynamically, but it verifies the
conjecture in the simplest non-trivial topological situation).
\end{theorem}

%% -------------------------------------------------------------------------
\section{Counterexample Attempts and Obstructions}
%% -------------------------------------------------------------------------

The difficulty of the Whitehead conjecture is illustrated by the fact that
there are many potential counterexample constructions that fail for
subtle reasons.

\begin{remark}
The simplest potential counterexample would be: an aspherical 2-complex $X$
and a subcomplex $Y \subseteq X$ with $\pi_2(Y) \neq 0$. For such $Y$,
$\pi_2(Y)$ would be a non-trivial $\Z[\pi_1(Y)]$-module. However, any
element $\alpha \in \pi_2(Y)$ maps to an element of $\pi_2(X) = 0$ (since
$X$ is aspherical), so $\alpha$ is already nullhomotopic in $X$.
The Whitehead conjecture says it is already nullhomotopic in $Y$.
\end{remark}

The key analytical difficulty is the following:

\begin{lemma}[Non-hereditary asphericity]
Asphericity is not hereditary under subspace inclusion in general.
For example, the torus $T^2 = K(\Z^2, 1)$ is aspherical, but if we
remove a 2-cell (to get the one-skeleton $S^1 \vee S^1$), we do get
an aspherical space. However, for more complicated examples, subcomplexes
can have non-trivial $\pi_2$.
\end{lemma}

\begin{example}
Consider a 2-complex $X$ that is the union of two aspherical subcomplexes
$A$ and $B$ along a common connected subcomplex $A \cap B$. By the
Mayer--Vietoris sequence in homotopy, $\pi_2(X) = 0$ does not immediately
imply $\pi_2(A) = 0$ or $\pi_2(B) = 0$ (the long exact sequence for
homotopy involves $\pi_1$-module structure). This is the source of difficulty.
\end{example}

%% -------------------------------------------------------------------------
\section{Open Problems}
%% -------------------------------------------------------------------------

\begin{conjecture}[Whitehead, Conjecture~\ref{conj:Whitehead}]
Every connected subcomplex of an aspherical 2-complex is aspherical.
\end{conjecture}

\begin{conjecture}[Eilenberg--Ganea, Conjecture~\ref{conj:EG}]
Every group $G$ with $\mathrm{cd}(G) = 2$ has a $K(G,1)$ of dimension 2.
\end{conjecture}

\begin{conjecture}[$D(2)$-Problem, Conjecture~\ref{conj:D2}]
Every finite 2-complex dominated by a 3-complex is homotopy equivalent
to a 2-complex.
\end{conjecture}

\begin{conjecture}[Asphericity of all one-relator groups]
\label{conj:one-relator-asph}
Every one-relator group $G = \langle S \mid r \rangle$ (without torsion)
has an aspherical presentation complex.
\end{conjecture}

\begin{conjecture}[Coherence Conjecture for Hyperbolic Groups]
Every one-relator group (equivalently, the fundamental group of every
compact surface-like 2-complex) is coherent.
\end{conjecture}

\begin{remark}
Conjectures~\ref{conj:Whitehead}, \ref{conj:EG}, and \ref{conj:D2} are
related: a counterexample to the Eilenberg--Ganea conjecture would give
a counterexample to Whitehead's conjecture; the $D(2)$-problem is equivalent
to Whitehead's conjecture for finite complexes (Johnson 2003).
A common approach to resolve all three would be a new construction of
aspherical complexes via CAT(0) geometry or hyperbolic group theory.
\end{remark}

%% -------------------------------------------------------------------------
\section{Conclusion}
%% -------------------------------------------------------------------------

The Whitehead asphericity conjecture, in spite of its elementary
formulation, touches on deep aspects of combinatorial group theory,
algebraic topology, and 4-manifold topology. The connection to the
$D(2)$-problem of Wall shows that it is not an isolated curiosity
but is interwoven with fundamental questions about the existence and
classification of 2-dimensional classifying spaces. Partial results
(Corson's theorem on diagrammatically reducible complexes, the
Magnus--Lyndon theorem for one-relator groups, Howie's result for
simply-connected subcomplexes) confirm the conjecture in important
special cases, but a general proof remains out of reach. The resolution
of even the special case of the Eilenberg--Ganea conjecture would be a
major advance.

\begin{thebibliography}{99}

\bibitem{Whitehead1941}
J.~H.~C.~Whitehead,
\textit{On adding relations to homotopy groups},
Ann.\ of Math.\ \textbf{42} (1941), 409--428.

\bibitem{Wall1965}
C.~T.~C.~Wall,
\textit{Finiteness conditions for CW-complexes},
Ann.\ of Math.\ \textbf{81} (1965), 56--69.

\bibitem{EilenbergGanea1957}
S.~Eilenberg, T.~Ganea,
\textit{On the Lusternik-Schnirelmann category of abstract groups},
Ann.\ of Math.\ \textbf{65} (1957), 517--518.

\bibitem{Magnus1930}
W.~Magnus,
\textit{Über diskontinuierliche Gruppen mit einer definierenden Relation
(der Freiheitssatz)},
J.\ Reine Angew.\ Math.\ \textbf{163} (1930), 141--165.

\bibitem{Lyndon1950}
R.~C.~Lyndon,
\textit{Cohomology theory of groups with a single defining relation},
Ann.\ of Math.\ \textbf{52} (1950), 650--665.

\bibitem{Cockcroft1954}
W.~H.~Cockcroft,
\textit{On two-dimensional aspherical complexes},
Proc.\ London Math.\ Soc.\ \textbf{4} (1954), 375--384.

\bibitem{Sieradski1976}
A.~J.~Sieradski,
\textit{A coloring invariant for topological graph theory},
J.\ Combin.\ Theory Ser.\ B \textbf{34} (1983), 241--246.

\bibitem{Corson1992}
J.~M.~Corson,
\textit{Conformally aspherical 2-complexes},
J.\ Pure Appl.\ Algebra \textbf{78} (1992), 11--22.

\bibitem{Scott1973}
G.~P.~Scott,
\textit{Finitely generated 3-manifold groups are finitely presented},
J.\ London Math.\ Soc.\ \textbf{6} (1973), 437--440.

\bibitem{BridsondelaHarpe1999}
M.~R.~Bridson, A.~de~la~Harpe,
\textit{Metric Spaces of Non-Positive Curvature},
Springer-Verlag, Berlin, 1999.

\bibitem{Howie1981}
J.~Howie,
\textit{Aspherical and acyclic 2-complexes},
J.\ London Math.\ Soc.\ \textbf{20} (1979), 549--558.

\bibitem{Johnson2003}
F.~E.~A.~Johnson,
\textit{Stable Modules and the $D(2)$-Problem},
London Mathematical Society Lecture Note Series \textbf{301},
Cambridge University Press, 2003.

\bibitem{Bestvina2002}
M.~Bestvina,
\textit{$\mathbb{R}$-trees in topology, geometry, and group theory},
in: \textit{Handbook of Geometric Topology}, Elsevier, 2002, 55--91.

\end{thebibliography}

\end{document}
